%% SAMPLE OUTPUT - DO NOT COPY INTO FINAL %%
%% Full LaTeX with filled preamble, short body, exhibit list %%
\documentclass[10pt,letterpaper,oneside]{extarticle}
\flushbottom
\usepackage{fontspec}
\usepackage{pdfpages}
\usepackage{graphicx}
\usepackage[left=1in,right=1in,top=1.05in,bottom=1.05in]{geometry}
\usepackage{tikz}
\usepackage{amssymb}
\usetikzlibrary{calc}
\usepackage{fancyhdr}
\usepackage{tocloft}
\usepackage{hyperref}
\usepackage{wasysym}
\usepackage{enumitem}
\usepackage{needspace}
\usepackage{microtype}
\usepackage{setspace}
\usepackage{array}
\usepackage{indentfirst}
\usepackage{booktabs}
\usepackage{footnote}
\usepackage{longtable}
\usepackage{xcolor}
\usepackage{titlesec}
\usepackage{ragged2e}

\titleformat{\section}
{\normalfont\fontsize{9}{10}\scshape\centering}{\thesection}{2em}{}
\titleformat{\subsection}
{\normalfont\fontsize{9}{10}\scshape\centering}{\thesubsection}{2em}{}
\titlespacing*{\section}{2pt}{5pt}{3pt}
\titlespacing*{\subsection}{2pt}{3pt}{2pt}
\renewcommand{\thesection}{\Roman{section}}
\renewcommand{\thesubsection}{\Alph{subsection}}

% ============================================================
% PLEADING LINE NUMBERS
% ============================================================
\newlength{\PleadingStep}
\setlength{\PleadingStep}{\dimexpr\textheight/28}
\makeatletter
\AtBeginDocument{\setlength{\baselineskip}{\PleadingStep}}
\makeatother
\newcommand{\PleadingTopInset}{5.5pt}
\newcommand{\PleadingNumsX}{0.95in}
\newcommand{\PleadingRuleXA}{0.95in}
\newcommand{\PleadingRuleXB}{0.975in}
\newcommand{\PleadingRuleWidth}{0.75pt}
\AddToHook{shipout/background}{%
	\begin{tikzpicture}[remember picture,overlay]
		\foreach \i in {1,...,28}{
			\pgfmathsetmacro{\k}{\i-1}
			\pgfmathsetlengthmacro{\yshift}{1in+\PleadingTopInset+(\k+0.5)*\PleadingStep}
			\node[anchor=east,font=\fontsize{10pt}{10pt}\selectfont] at
			([xshift=\PleadingNumsX,yshift=-\yshift]current page.north west){\i};
		}
		\pgfmathsetlengthmacro{\yTop}{1in+\PleadingTopInset}
		\pgfmathsetlengthmacro{\yBot}{1in+\PleadingTopInset+\textheight}
		\draw[line width=\PleadingRuleWidth]
		([xshift=\PleadingRuleXA,yshift=-\yTop]current page.north west)--
		([xshift=\PleadingRuleXA,yshift=-\yBot]current page.north west);
		\draw[line width=\PleadingRuleWidth]
		([xshift=\PleadingRuleXB,yshift=-\yTop]current page.north west)--
		([xshift=\PleadingRuleXB,yshift=-\yBot]current page.north west);
	\end{tikzpicture}%
}

% ============================================================
% HEADER / FOOTER
% ============================================================
\pagestyle{fancy}
\fancyhf{}
\renewcommand{\headrulewidth}{0pt}
\renewcommand{\footrulewidth}{0.4pt}
\fancyfoot[L]{\fontsize{8}{10}\selectfont
	\textbf{\textsc{DEFENDANT'S AMENDED ANSWER; AFFIRMATIVE DEFENSES;
			RENEWED MOTION TO DISMISS}} --- Pg.~\thepage}
\setlength{\headheight}{13.6pt}

% ============================================================
% FORMATTING
% ============================================================
\setlength{\parindent}{0.5in}
\setlength{\parskip}{.05pt}
\hyphenpenalty=5000
\tolerance=1000
\emergencystretch=3em

% ============================================================
% CUSTOM COMMANDS
% ============================================================
\newcommand{\CenteredUnderlined}[1]{%
	\begin{center}\underline{\textbf{#1}}\end{center}}
\newcommand{\HRule}{\noindent\rule{\linewidth}{0.5pt}}
\newcommand{\CB}{\Square\;}
\newcommand{\CaptionBlock}[3]{%
	\noindent
	\begin{minipage}[t]{3.2in}
		\vspace{12pt}%
		\begin{tikzpicture}
			\node[inner sep=0pt,text width=3.2in,anchor=north west](T){#3};
			\draw[line width=0.5pt](T.south west)--(T.south east);
			\draw[line width=0.5pt](T.south east)--(T.north east);
		\end{tikzpicture}
	\end{minipage}%
	\hfill
	\begin{minipage}[t]{3.0in}
		\vspace{2pt}%
		\RaggedRight Case No.: #1\\[10pt]
		\raggedleft\small\textbf{#2}
	\end{minipage}%
}

% ===========================================================
\begin{document}
% ============================================================
% DOCUMENT 1 OF 4 — AMENDED ANSWER
% ============================================================
\begin{center}
	{\normalsize\textbf{IN THE CIRCUIT COURT OF THE STATE OF OREGON}}\\
		{\normalsize\textbf{FOR THE COUNTY OF WASHINGTON}}\\
\end{center}
\CaptionBlock{\textbf{26LT03815}}{%
	\begin{center}
		\fontsize{9}{9}\selectfont
		\textbf{DEFENDANT'S AMENDED ANSWER;\\
			AFFIRMATIVE DEFENSES; AND\\
			RENEWED MOTION TO DISMISS FOR\\
			LACK OF SUBJECT MATTER JURISDICTION\\
			AND LACK OF STANDING}%
	\end{center}
}{%
	JIM WONG,\\[4pt]
	\hspace{0.3in}Plaintiff,\\[8pt]
	vs.\\[8pt]
	ASHLE' PENN,\\[8pt]
	\hspace{0.3in}Defendant, \textit{In Propria Persona}.%
}

\begin{center}
	\fontsize{9}{9}\selectfont
	\textbf{SPECIAL APPEARANCE ---
		DEFENDANT DOES NOT WAIVE JURISDICTIONAL DEFENSES}\\[2pt]
	\textbf{ORCP 21 A(1)(a), (b), (d), (e), (f), (g);
		ORCP 21 G(1), G(2), G(4)}\\[3pt]
	\textbf{AMENDMENT AS OF RIGHT PURSUANT TO ORCP 23~A}\\[3pt]
	\textit{No responsive pleading has been served by Plaintiff.
		Defendant amends as of right. This Amended Answer
		supersedes and replaces the original Answer filed
		March~2, 2026, and is construed to have been filed
		on that date. ORCP 23~A; ORCP 25.}
\end{center}
\indent COMES NOW Defendant Ashle' Penn (Defendant''),
appearing specially and without waiving any jurisdictional
objection previously raised or preserved, and for her Amended
Answer states as follows.
% ============================================================
\section{\textbf{PRELIMINARY MATTERS}}
% ============================================================
\textbf{Amendment as of Right.}
No responsive pleading has been served by Plaintiff. This
Amended Answer is filed as of right under ORCP 23~A. It
supersedes and replaces the original Answer filed March~2,
2026, and relates back to that date for all purposes.
\footnote{ORCP 23~A}; \footnote{ORCP 25}. This filing does not constitute a general
appearance and does not waive any jurisdictional defense
previously raised.\\
\textbf{Filing Under Compulsion.}
This Amended Answer is filed solely under compulsion of an
oral deadline imposed during the March~2, 2026 telephonic
hearing. That deadline is objected to on three independent
grounds. \textit{First}, under UTCR~21.080(2), the deadline
for electronically filed documents is 11:59:59~p.m.; an oral
instruction cannot supersede a statewide rule.
\textit{Second}, the deadline was imposed after Plaintiff had
dropped from the call --- the case had been called, Plaintiff
was no longer present, and all substantive action was taken
with only Defendant on the line, in violation of the
prohibition on ex parte proceedings. \textit{Third},
Plaintiff communicated the trial date to Defendant by text
message at 10:19~a.m.\ --- 27 minutes before the court
called Defendant at 10:46~a.m.\ --- stating: "I spoke with
the court regarding your motion to file documents, and the
court has scheduled a trial date for June 5, 2026.''
\footnote{\textit{See} Exhibit~M}. This sequence establishes that
Plaintiff had ex parte contact with the court concerning
Defendant's pending motion before Defendant was contacted,
and that the trial date was set as a result of that
communication. Oregon Code of Judicial Conduct Rule~2.9(A)
prohibits a judge from initiating, permitting, or considering
ex parte communications concerning a pending matter. This
filing does not constitute, and shall not be construed as,
a waiver of any jurisdictional defense previously raised,
including those in Defendant's Motions to Dismiss filed
March~2 and March~5, 2026, each incorporated herein by
reference. \footnote{ORCP 21~G(1), G(4)}.
\textbf{Fee Objection.}
Under ORS~\S105.130(3), the \$88 fee applies only upon a
demand for trial. This filing is not a trial demand. Any fee
assessed is paid under protest and without waiver.
\textbf{Trial Date Objection.}
When this case was called on March~2, 2026, Plaintiff was not
present on the line. The court itself confirmed: We lost
the plaintiff.'' ORS~\S105.137(3) mandates that the court
\textit{shall} dismiss the action or enter judgment for the
defendant upon the plaintiff's failure to appear. That
provision is non-discretionary. Defendant's presence was
special and limited to jurisdictional objections only and
does not constitute a general appearance. The court was not
authorized to set a trial date over a pending, unresolved
jurisdictional motion\footnote{ORCP~21~G(4)}; \footnote{ORS~\S105.130(4)}.
The docket reflects the trial date order as entered by
Judge Meisner on March~2, 2026; the trial is assigned to
Judge Knaupp on June~5, 2026. Defendant's Motion to Dismiss
has not been ruled upon and remains pending before this
Court.
\textbf{HAP Payments --- Plaintiff's Admission.}
Plaintiff's own written communications confirm that he
received HAP payments from WCHA throughout the tenancy.
\textit{See} Exhibit~C. Those payments are recorded as
\$0.00 received in every row of Plaintiff's ledger.
\textit{See} Exhibit~B. The discrepancy between payments
admitted and payments credited is not a bookkeeping error;
it is the mechanism by which the manufactured arrearage was
constructed.
% ============================================================
\section{\textbf{PRESERVATION OF JURISDICTIONAL DEFENSES}}
% ============================================================
Pursuant to ORCP 21~A(1) and ORCP 21~G, Defendant expressly
preserves:
\begin{enumerate}[label=(\textbf{\alph*)},leftmargin=.5in,nosep]
	\item Lack of subject matter jurisdiction ---
	ORCP 21~A(1)(a); ORCP 21~G(4) (non-waivable; must be
	raised at any time);
	\item Lack of jurisdiction over the person ---
	ORCP 21~A(1)(b); ORCP 21~G(1);
	\item Insufficiency of summons and service of process ---
	ORCP 21~A(1)(e); ORCP 21~G(1);
	\item Plaintiff lacks legal capacity and is not the real
	party in interest --- ORCP 21~A(1)(d), (f);
	ORCP 21~G(2);
	\item Failure to state a claim upon which relief can be
	granted --- ORCP 21~A(1)(g); and
	\item Failure to join an indispensable party ---
	ORCP 21~A(1)(g); ORCP 26~A.
\end{enumerate}
\indent No defense is waived by joinder with other defenses.
ORCP 21~A(2)(a).
% ============================================================
\section{\textbf{RENEWED MOTION TO DISMISS}}
% ============================================================
\subsection{\textbf{Defective Notice --- No Subject Matter
		Jurisdiction
		(ORS~\S90.394(3); \textit{Hickey v.\ Scott},
		370 Or 97 (2022))}}
The notice of termination states \$5,918.00 as the amount
of rent owed. That figure is derived from a ledger that
embeds \$2,220.00 in HOA fines and \$1,045.00 in late fees
--- neither of which constitutes rent'' under
ORS~\S90.100(38) --- and applies the \$4,494.00 HRSN
payment to November 2025 only, in violation of the mandatory
chronological application order of ORS~\S90.220(9). The
ledger also includes charges for November 2025 through
January 2026, three months in which Defendant's recertified
tenant portion was \$0.00. \textit{See} Exhibits~B and~G.\\
\indent \textbf{Late fees are independently prohibited from
	appearing in a termination notice.} ORS~\S90.260(4)
provides that a landlord may not include a late charge in a
notice of termination for nonpayment of rent. \indent Plaintiff's own ledger, read against the statutes
governing what may lawfully appear in a termination notice,
yields a maximum cognizable balance of \$1,345.00---and
that figure assumes, without conceding, that every
remaining charge is valid.

\begin{center}
	\small
	\begin{tabular}{@{}p{4.2in}rp{1.3in}@{}}
		\toprule
		\textbf{Category} & \textbf{Amount} &
		\textbf{Authority} \\
		\midrule
		Tenant-portion rent, Feb.~2025--Oct.~2025
		(\$607 $+$ (8 $\times$ \$654))
		--- HAP payments recorded as \$0.00
		received despite Plaintiff's admission
		of receipt
		& \$5,839.00 & Exhibit B \\[4pt]
		\textit{Less:} HRSN \$4,494.00 applied
		chronologically from Feb.~2025;
		exhausts mid-Aug.~2025
		(\$617 of \$654; \$37 remainder unpaid)
		& (\$4,494.00) & ORS \S90.220(9) \\[4pt]
		\midrule
		\textbf{Maximum cognizable balance per
			Plaintiff's ledger, prohibited charges
			excluded} (\$37 $+$ \$654 $+$ \$654)
		& \textbf{\$1,345.00} & \\
		\bottomrule
	\end{tabular}
\end{center}

\indent Plaintiff did not notice \$1,345.00. Plaintiff
noticed \$5,918.00---an overstatement of \$4,573.00.
That overstatement is not rounding error or bookkeeping
imprecision. It is the product of four categories of
charges the Legislature has expressly prohibited from
inclusion in a termination notice. Each is itemized below.

\begin{center}
	\small
	\begin{tabular}{@{}p{3.8in}rp{1.5in}@{}}
		\toprule
		\textbf{Prohibited Charge} & \textbf{Amount} &
		\textbf{Authority} \\
		\midrule
		HOA fines/assessments (8 months,
		Jul.~2024--Apr.~2025) --- not rent under
		ORS~\S90.100(41)(a); pass-through
		prohibited under ORS~\S90.302(1)
		& \$2,220.00 &
		ORS \S\S90.100(41)(a), 90.302(1) \\[4pt]
		Late fees (\$95 $\times$ 11 months) ---
		prohibited in termination notice
		& \$1,045.00 & ORS \S90.260(4) \\[4pt]
		Nov.~2025--Jan.~2026 charges at \$0.00
		tenant portion post-recertification
		(3 $\times$ \$654)
		& \$1,962.00 & Exhibits G, H \\[4pt]
		HRSN \$4,494.00 misapplied to
		Nov.~2025 only instead of
		chronologically from Feb.~2025
		& \$4,494.00 & ORS \S90.220(9) \\[4pt]
		\midrule
		\textbf{Total overstatement}
		& \textbf{\$4,573.00} & \\
		\bottomrule
	\end{tabular}
\end{center}
\indent To the extent the Court reaches the question of
rent owed notwithstanding the void notice---which Defendant
contests---the corrected balance derived from Plaintiff's
own ledger is \textbf{\$1,345.00}. To the extent the Court
reaches that corrected balance, it is fully offset and
negated by Defendant's affirmative defenses and
counterclaims, which independently entitle Defendant to a
net judgment in her favor. \footnote{\textit{See} Counterclaims,
	\textit{infra}.}\\
\indent The notice overstates the amount owed by \$4,573.00
--- approximately 77\% --- and is void under
\textit{Hickey}. This Court lacks subject matter
jurisdiction as a result. Plaintiff's January~21, 2026 email transmitting the ledger states: the HOA kindly agreed to temporarily
	suspend additional fine HOA account.'' \textit{See}
	Exhibit~C. This is a direct admission that HOA fine
	pass-throughs were within Plaintiff's discretion and
	control, not mandatory charges, and that Plaintiff
	knew they were not rent at the time he included them
	in the noticed balance.}
\subsection{\textbf{Wrong Plaintiff; Failure to State a
		Claim; Indispensable Party Not Joined
		(ORCP 21~A(1)(d), (f), (g); ORS~\S105.130(4);
		ORS~\S63.165; ORCP 26~A)}}
The lease identifies the lessor as \textbf{JWP WONG REAL
	ESTATE}, an Oregon limited liability company.
\textit{See} Exhibit~D. HAP payments from WCHA and the
\$4,494.00 HRSN benefit were each paid directly to JWP WONG
REAL ESTATE --- not to Jim Wong individually.
\textit{See} Exhibit~G. Jim Wong individually is not the
real party in interest and lacks legal capacity to maintain
this action. ORS~\S63.165 precludes an LLC member from
asserting the LLC's legal rights in his own name absent
express authorization. Jim Wong has produced no operating
agreement, board resolution, or agency authorization. The
LLC is an indispensable party not joined. ORCP 26~A;
ORCP 21~A(1)(g).\\
\indent \textbf{Failure to state a claim arises on two
	independent grounds.} \textit{First}, a plaintiff who is
not the real party in interest has no cognizable legal right
to possession and therefore cannot state a claim for
restitution of the premises. ORS~\S105.130(4) requires that
the action be brought by the landlord'' --- the party with
whom the tenant has a rental agreement, which is JWP WONG
REAL ESTATE. \textit{Second}, a complaint premised on a
termination notice that is void under ORS~\S90.394(3) and
\textit{Hickey} cannot state a cognizable claim for
possession; the void notice is the jurisdictional predicate
for the entire action, and without it there is no claim.
ORCP 21~A(1)(g).
\indent Defendant requests dismissal without prejudice as
to the real-party defect, or substitution of JWP WONG REAL
ESTATE with proof of agency authority before further
proceedings. The real-party defect runs against Plaintiff's
claim only and is not available to Plaintiff as a shield
against Defendant's compulsory counterclaims.
ORS~\S90.370(1)(a); ORCP 22~A.
\subsection{\textbf{Additional Grounds for Dismissal}}
\textbf{Defective Service.}
Service did not comply with ORCP~7 and ORS~\S105.135,
depriving this Court of personal jurisdiction over
Defendant. ORCP 21~A(1)(e); ORCP 21~G(1). This objection
was raised on the record March~2, 2026 and is preserved.\\
\textbf{Failure to Cooperate with Rental Assistance.}
ORS~\S90.395 independently requires dismissal where the
landlord failed to cooperate with the tenant's rental
assistance program. Plaintiff accepted both the HRSN
payment and monthly HAP payments from WCHA, then
misapplied those funds to unlawful charges in violation of
ORS~\S90.220(9)'s mandatory payment application order, and
filed this action claiming a balance that does not reflect
payments received. Accepting rental assistance and
misapplying it to manufactured charges is not cooperation
within the meaning of ORS~\S90.395. Plaintiff's conduct
also constitutes a breach of the HAP contract between JWP
WONG REAL ESTATE and WCHA under 24~CFR~\S982.451, and may
independently constitute an owner breach of contract under
24~CFR~\S982.453. Defendant will provide notice to WCHA of
these violations.\\
\textbf{Plaintiff's Non-Appearance; ORS~\S105.137(2).}
When this case was called on March~2, 2026, Plaintiff was
not present on the line. The court confirmed: "We lost the
Plaintiff.'' The court then proceeded to take substantive
action --- setting a trial date, imposing an answer
deadline, and overruling Defendant's jurisdictional
objection --- with only Defendant on the line.
ORS~\S105.137(2) is mandatory: upon Plaintiff's failure to
appear, the court \textit{shall} dismiss or enter judgment
for Defendant. The court's failure to apply
ORS~\S105.137(2) is preserved error. Plaintiff's 10:22~a.m.\
text message --- sent 27 minutes before the court called
Defendant --- confirms that Plaintiff had ex parte contact
with the court regarding Defendant's pending motion before
Defendant was reached. \textit{See} Exhibit~M. Defendant's
objection to the trial date was stated on the record,
overruled, and is preserved for appellate review. No trial
date should have been set while Defendant's jurisdictional
motion remained pending and unresolved and after Plaintiff was not present during the hearing. ORCP~21~G(4).
% ============================================================
\section{\textbf{ANSWER AND DENIAL OF ALLEGATIONS}}
% ============================================================
Without waiving any jurisdictional defense, Defendant
responds to Plaintiff's Complaint as follows:
\begin{enumerate}[label=\arabic*.,leftmargin=.5in,nosep]
	\item Defendant \textbf{Denies} that Plaintiff Jim Wong
	individually is entitled to possession of 17548 NW
	Springville Rd \#9, Portland, OR 97229.
	\item Defendant \textbf{Denies} that the landlord of
	record is Jim Wong. The landlord of record is JWP WONG
	REAL ESTATE, LLC.
	\item Defendant \textbf{Denies} that she owes the
	amount stated in the notice of termination or Complaint.
	\item Defendant \textbf{Denies} that the notice of
	termination was valid, legally sufficient, or properly
	served.
	\item Defendant \textbf{Denies} each remaining
	allegation not specifically admitted herein.
	\item Defendant \textbf{Admits} that Plaintiff Jim Wong
	is the titled owner of the unit.
\end{enumerate}
% ============================================================
\section{\textbf{AFFIRMATIVE DEFENSES}}
% ============================================================
\indent This Answer is filed under court compulsion. This
court-compelled filing does not constitute a voluntary
general appearance. The Court is obligated to resolve the
jurisdictional, standing and capacity questions first. All affirmative defenses and
counterclaims are preserved for any subsequent proceeding
in which jurisdiction is properly established.
\subsection{\textbf{First Affirmative Defense ---
		Defective Notice
		(ORS~\S90.394(3); \textit{Hickey v.\ Scott},
		370 Or 97 (2022))}}
The notice overstates the amount of rent owed by \$4,573.00.
With payments applied consistent with statutory requirements, any alleged balance would amount to \$1,345.00. Under
\textit{Hickey v.\ Scott}, 370 Or 97 (2022), any
overstatement of the amount of rent owed voids the notice
and deprives this Court of subject matter jurisdiction.
\textit{See} Section~III(A), \textit{supra}.
\subsection{\textbf{Second Affirmative Defense ---
		Misapplication of Payments; Unlawful Late Fees
		(ORS~\S90.220(9); ORS~\S90.260(4);
		ORS~\S90.395)}}
Plaintiff received HAP payments from WCHA each month and
the \$4,494.00 HRSN benefit, yet applied those funds to HOA
fees and late charges ahead of rent, in violation of
ORS~\S90.220(9)'s mandatory payment application order.
\textit{See} Exhibits~B and~G. ORS~\S90.260(4) independently
prohibits the inclusion of late charges in a termination
notice for nonpayment of rent. The \$1,045.00 in late fees
in Plaintiff's noticed balance violates ORS~\S90.260(4)
regardless of whether the underlying late fee provision is
otherwise valid. Correctly applied, the balance is
\$1,345.00. Plaintiff's failure to cooperate with
Defendant's rental assistance programs independently
requires dismissal under ORS~\S90.395.
\subsection{\textbf{Third Affirmative Defense ---
		Zero Tenant Portion Following Recertification}}
Defendant recertified her income with WCHA in November 2025,
reporting zero income, reducing her tenant portion to
\$0.00 for November 2025, December 2025, and January 2026.
\textit{See} Exhibits~G and~H. Plaintiff's inclusion of
\$1,962.00 for those three months independently overstates
the noticed balance and voids the notice under
\textit{Hickey}.
\subsection{\textbf{Fourth Affirmative Defense ---
		Failure to Maintain Habitable Premises
		(ORS~\S90.320; ORS~\S90.360(2); ORS~\S90.370;
		ORS~\S90.380)}}
Plaintiff failed to maintain the premises in habitable
condition as required by ORS~\S90.320. Defendant provided
written notice of each of the following defects, each of
which remained unresolved as of this filing:
\begin{enumerate}[label=(\alph*),leftmargin=.5in,
	itemsep=4pt]
	\item \textbf{Drain fly infestation from plumbing
		defect.} Reported November~16, 2023; unresolved as of
	this filing --- approximately \textbf{28 months}.
	ORS~\S90.320(1)(a)(b)(e)(f) (plumbing facilities in good working
	order). \textit{See} Exhibit~C.
	\item \textbf{Missing window screens.} Requested
	July~2023; unresolved as of this filing ---
	approximately \textbf{32 months}.
	ORS~\S90.320(1)(a) (windows in good repair).
	Plaintiff's 30-month failure to install window screens
	directly caused the HOA window-covering noncompliance
	notices that Plaintiff then forwarded to WCHA.
	\textit{See} Exhibits~C, I--L.
	\item \textbf{Leaking kitchen sink.} Reported
	June~26, 2024; unresolved as of this filing ---
	approximately \textbf{21 months}.
	ORS~\S90.320(1)(c). \textit{See} Exhibit~C.
	\item \textbf{Water intrusion into live electrical
		outlet.} Reported June~26, 2024; unresolved as of this
	filing --- approximately \textbf{21 months}.
	ORS~\S90.320(1)(e) (electrical systems in good working
	order). This defect constitutes an immediate and
	ongoing safety hazard. \textit{See} Exhibit~C.
	\item \textbf{Leaking master toilet; standing water;
		mold.} Reported November~17, 2025; unresolved as of
	this filing --- approximately \textbf{4 months}.
	ORS~\S90.320(1)(b) (plumbing; elimination of
	moisture and mold) (f). \textit{See} Exhibit~C.
	\item \textbf{Non-functioning stove burner; inadequate
		hot water.} Reported November~2025--January~2026;
	unresolved as of this filing --- approximately
	\textbf{4 months}. ORS~\S90.320(c)(A) (heating and
	cooking facilities in good working order).
	\textit{See} Exhibit~C.
	\item \textbf{Intermittent power outages;
		non-functioning outlets.} Reported January~6, 2026;
	unresolved as of this filing --- approximately
	\textbf{2 months}. ORS~\S90.320(1)(e).
	\textit{See} Exhibit~C.
	\item \textbf{Failure to provide replacement mail key.}
	Three written requests; unresolved as of this filing.
	ORS~\S90.320(1) (premises fit for human occupation).
	This failure deprived Defendant of legal, court, and
	government correspondence throughout this proceeding.
	\textit{See} Exhibit~C.
\end{enumerate}
\indent ORS~\S90.360(2) provides that where a landlord
fails to remedy a habitability defect after written notice,
the tenant may recover the diminution in rental value of
the dwelling for the period of the defect. A conservative
10\% diminution of \$2,535.00/month applied to the full
32-month tenancy yields \$8,112.00 in rent abatement. The
10\% figure is conservative; the court may find a higher
percentage warranted given the number, duration, and
severity of the defects, particularly the 28-month drain
fly infestation and the live electrical hazard. Each item
above was noticed in writing. \textit{See} Exhibit~C.\\
\indent \textbf{Renting dwelling known to be unlawful to
	occupy (ORS~\S90.380).} Plaintiff knew or reasonably should
have known that the unit was unlawful to occupy given the
water intrusion into a live electrical outlet (reported
June~26, 2024; unresolved 21 months), the 28-month drain
fly infestation originating from a plumbing defect, and the
standing water and mold from the leaking master toilet.
ORS~\S90.380 provides that where a landlord rents a
dwelling that the landlord knows or reasonably should know
is unlawful to occupy, the tenant may recover the greater
of two months' periodic rent (\$5,070.00) or twice the
actual damages sustained. The floor of this remedy is
\$5,070.00; actual damages are to be proven at trial.\\
\indent The causal chain is direct: Plaintiff created the
conditions generating the HOA fines by failing to install
window screens for 32 months, charged those fines to
Defendant in violation of ORS~\S90.302(1), applied
Defendant's payments to those unlawful charges first in
violation of ORS~\S90.220(9), and then filed this eviction
claiming unpaid rent. Each step is independently
%% PICKS UP from: "Each step is independently"
unlawful; together they are a manufactured arrearage.
ORS~\S90.370; ORS~\S90.385.
\subsection{\textbf{Fifth Affirmative Defense ---
		Retaliation; Selective Enforcement; Discrimination
		(ORS~\S90.385)}}
This eviction follows Defendant's assertion of habitability
rights and constitutes retaliation under ORS~\S90.385 on
four independent bases:
\textbf{(1) Selective HCV Disclosure.}
In November 2025, Plaintiff selectively forwarded an HOA
window-covering noncompliance notice to WCHA at
hcvinterim@washingtoncountyor.gov --- a violation caused
entirely by Plaintiff's own 32-month failure to install
window screens --- while simultaneously concealing from
WCHA eleven documented GHQS violations under
24~CFR~\S982.401, any one of which would have required an
HQS inspection and exposed Plaintiff to suspension of HAP
payments. \textit{See} Exhibits~C, I--L. This selective,
bad-faith disclosure was designed to threaten Defendant's
Section~8 voucher while shielding Plaintiff from regulatory
consequences.
\textbf{(2) Selective Enforcement of HOA Rules --- Grill.}
The rental agreement between Defendant and JWP WONG REAL
ESTATE expressly permits Defendant to maintain an outdoor
grill on the premises. \textit{See} Exhibit~D. The patio
is Defendant's exclusive-use limited common element under
CC\&R Sections~5.1 and~6.1. \textit{See} Exhibit~E.
Plaintiff enforced an HOA demand that Defendant remove her
grill --- in direct contradiction of a written lease term
--- while other residents of the complex are permitted to
maintain grills and outdoor equipment without comparable
enforcement action. \textit{See} Exhibit~N. This selective
enforcement, directed at Defendant and not at similarly
situated residents, constitutes: (a)~breach of the rental
agreement; (b)~retaliation under ORS~\S90.385;
(c)~a violation of Defendant's quiet enjoyment rights under
ORS~\S90.360(1); and (d)~a basis for lease termination
without penalty under ORS~\S90.360(1). A landlord cannot
simultaneously grant a contractual right in the lease and
then enforce an HOA demand to extinguish it. Defendant
reserves the right to assert a fair housing claim based on
the pattern of selective enforcement if discovery reveals
that enforcement has been directed disproportionately at
protected-class tenants.
\textbf{(3) Bad-Faith Lease Coercion.}
Plaintiff engaged in bad-faith lease coercion in writing on
or about January~25, 2025. \textit{See} Exhibit~C.
\textbf{(4) Prior Defective FED Judgment;
	Constructive Entrapment.}
Plaintiff obtained a prior FED judgment in Washington County
Circuit Court through service directed to a neighboring
address at which Defendant has never resided. That judgment
clouds Defendant's rental history, prevents her from
securing alternative housing, and effectively traps her in
this tenancy. Plaintiff has repeatedly refused to stipulate
to vacatur despite Defendant's written requests ---
deliberately maintaining a cloud he knows to be infirm ---
while simultaneously filing this new action. A landlord who
traps a tenant in a substandard unit by refusing to vacate
a void judgment, and then files a second eviction for the
same tenancy, does not come to this Court with clean hands.
\subsection{\textbf{Sixth Affirmative Defense ---
		Unlawful Vehicle Tow; Conversion
		(ORS~\S90.325; ORS~\S90.365;
		ORS~\S98.854(2); ORS~\S12.080)}}
On or about March~6, 2024, Defendant's Dodge Journey was
towed from the Courtyards at Springville by Sergeant's
Towing at AMS direction. Plaintiff owed Defendant a
non-delegable duty under ORS~\S90.325 to ensure that
third-party enforcement actions directed at Defendant's
tenancy --- including towing --- complied with her statutory
and contractual rights. Plaintiff breached that duty in the
following independently sufficient respects, each documented
by Plaintiff's own written admissions
(\textit{see} Exhibit~C):
\textbf{Duty.} As the unit owner and landlord in privity
with Defendant, Plaintiff bore a non-delegable duty of
reasonable care under ORS~\S90.325 to supervise and control
the conduct of AMS and any towing company acting at AMS
direction with respect to Defendant's property.\\
\textbf{Breach.}
\begin{enumerate}[label=(\alph*),leftmargin=.5in,
	itemsep=4pt]
	\item \textbf{No parking pass.} From move-in in
	July~2023, Defendant registered her vehicles with
	Plaintiff in writing. Despite repeated requests over
	more than eight months, Plaintiff never issued a
	parking pass. Plaintiff's own March~11, 2024 email
	requesting Defendant fill out a vehicle registration
	form confirms no valid pass had been issued as of
	the date of the tow. \textit{See} Exhibit~C.
	\item \textbf{No valid tow authorization.} Under
	ORS~\S98.854(2), a tower may not remove a vehicle
	without signed authorization from the property owner
	or owner's agent at the time of the tow. No
	authorization from Plaintiff --- the party in privity
	with Defendant --- was obtained.
	\item \textbf{No pre-tow warning.} Plaintiff
	acknowledged in writing on March~17--18, 2024 that a
	warning was required before towing and that
	Defendant's photographs showed her vehicle was not in
	the fire lane. \textit{See} Exhibit~C.
	\item \textbf{Shifting justifications.} The stated
	basis for the tow shifted at least three times (fire
	lane; blocking pedestrian access; failure to display
	tags), none supported by the date- and
	time-stamped photographs required by
	ORS~\S98.853(2). \textit{See} Exhibit~N.
	\item \textbf{Abdication.} Plaintiff told Defendant
	in writing: "
	there is nothing I could do,'' and
	took no action to remedy the tow or recover the
	vehicle. \textit{See} Exhibit~C (March~17, 2024).
\end{enumerate}
\textbf{Causation and Damages.} Plaintiff's breach of his
non-delegable duty directly caused the loss of Defendant's
Dodge Journey and all personal property contained therein,
including an iPad, prescription eyeglasses, prescription
medications, legal and financial paperwork, clothing, a
child's car seat, jewelry, and work equipment. The
outstanding payoff balance on the vehicle --- reported to
Defendant's credit --- is \$14,039.88. \textit{See}
Exhibit~P. The loss of the vehicle caused further
consequential damages including credit impairment, loss of
transportation, impaired employment mobility, and emotional
distress, to be proven at trial.\\
\textbf{Statute of Limitations.} This claim is governed by
ORS~\S12.080, which provides a six-year limitations period
for claims arising from the conversion of personal property.
The tow occurred March~6, 2024; the limitations period
expires March~6, 2030. This filing is timely.\\
\textbf{Forthcoming Motion to Compel.} Defendant will file
a Motion to Compel production of all towing authorization
records, communications between Plaintiff and AMS or
Sergeant's Towing, and any HOA enforcement records relating
to the March~6, 2024 tow.
\subsection{\textbf{Seventh Affirmative Defense ---
		Unlawful Entry into Exclusive-Use Space;
		Breach of Quiet Enjoyment;
		Unlawful Fee Pass-Through
		(ORS~\S90.322(1)(f), (8); ORS~\S90.360(1);
		ORS~\S90.302(1), (3); ORS~\S90.325;
		CC\&R \S\S5.1, 6.1)}}
Plaintiff permitted, authorized, or failed to prevent the
HOA and AMS from repeatedly entering Defendant's
exclusive-use patio, backyard, and garage --- limited common
elements for Defendant's exclusive use under CC\&R
Sections~5.1 and~6.1 (\textit{see} Exhibit~E) --- without
providing the 24-hour notice required by
ORS~\S90.322(1)(f) and without Defendant's consent.
Plaintiff's own July~30, 2025 text message confirms the
mechanism: They will send compliance officer to check on
whether they have been addressed and accumulate the fines
each month.'' \textit{See} Exhibit~C. Plaintiff's ledger
reflects eight discrete months of HOA fine assessments
(July~2024--April~2025), each corresponding to a compliance
inspection of Defendant's exclusive-use space without
required notice, establishing a minimum of eight violations.
\textit{See} Exhibit~B. A pre-ledger entry on
September~25, 2023 establishes at least one additional
violation.\\
\indent \textbf{Breach of quiet enjoyment.} Each
unauthorized entry into Defendant's exclusive-use patio,
backyard, and garage independently constitutes a breach of
Defendant's right to quiet enjoyment under
ORS~\S90.360(1), entitling Defendant to terminate the
rental agreement without penalty and to recover actual
damages in addition to the statutory remedy under
ORS~\S90.322(8). The CC\&Rs govern the relationship between
unit owners and the HOA --- not the relationship between
Plaintiff and Defendant. Defendant's rights are governed
exclusively by the ORLTA. The HOA's authority to conduct
compliance inspections of common areas does not extend to
Defendant's exclusive-use limited common element without
first complying with ORS~\S90.322(1)(f). Plaintiff, as unit
owner and landlord, is responsible under ORS~\S90.325 for
ensuring that HOA enforcement directed at Defendant's
exclusive-use space complies with Defendant's statutory
rights.\\
\indent \textbf{Unlawful fee pass-through.}
ORS~\S90.302(1) prohibits a landlord from charging,
collecting, or attempting to collect any fee not authorized
by statute or rental agreement. HOA assessments run against
unit owners only under Bylaws Sections~5.3(a), 5.6, and~5.7
(\textit{see} Exhibit~F); Plaintiff has no authority to
pass them through to Defendant. Any lease clause purporting
to authorize such pass-throughs is void under ORS~\S90.135
as an attempted waiver of Defendant's ORLTA rights. Under
ORS~\S90.302(3), a landlord who charges, collects, or
attempts to collect a fee not authorized by statute or
rental agreement is liable for two months' periodic rent
(\$5,070.00) plus actual damages (\$2,220.00 confirmed on
Plaintiff's own ledger) per claim. The demand itself is the
statutory violation --- payment is not required to trigger
the remedy.\\
\indent \textbf{Statute of Limitations.} All HOA entry and
fee claims are timely. ORS~\S12.125 provides a one-year
limitations period for claims arising under ORS Chapter~90
or a rental agreement. Entries and demands within one year
of this filing are timely on their face. Entries and demands
beyond one year are timely under the continuing tort
doctrine: each unauthorized entry and each unlawful fee
demand was a new and independent violation, and the claim
accrued anew with each occurrence.\\
\indent \textbf{Forthcoming Motion to Compel.} Defendant
will file a Motion to Compel production of all HOA
compliance inspection records, fining records, and
communications between Plaintiff, AMS, and the HOA relating
to Unit~9 from July~2023 to present.
\subsection{\textbf{Eighth Affirmative Defense ---
		Unclean Hands; Bad Faith;
		Constructive Entrapment}}
Plaintiff seeks equitable relief --- a writ of possession
--- while having engaged in a sustained course of
inequitable conduct that is the direct cause of the
conditions he now complains of:
\begin{enumerate}[label=(\arabic*),leftmargin=.5in,
noitemsep]
	\item Plaintiff obtained a prior FED judgment in
	Washington County Circuit Court through service
	directed to a neighboring address, producing a void
	judgment that has clouded Defendant's rental history
	and prevented her from securing alternative housing;
	\item Plaintiff has repeatedly refused to stipulate to
	vacatur of that void judgment, deliberately maintaining
	a cloud he knows to be infirm;
	\item Plaintiff filed the present action without
	enforcing the prior judgment --- a functional
	concession that the prior judgment cannot stand ---
	while simultaneously refusing to vacate it, trapping
	Defendant in this tenancy;
	\item Plaintiff seeks possession of a unit he has
	allowed to fall into uninhabitable condition for over
	32 months while Defendant had no viable ability to
	leave; and
	\item Plaintiff created the very rent arrearage he
	relies upon by permitting unlawful HOA entries,
	passing through the resulting fines in violation of
	ORS~\S90.302(1), and misapplying Defendant's payments
	in violation of ORS~\S90.220(9).
\end{enumerate}
\vspace{-1em}
\indent A party who manufactures the conditions making a
tenant's continued occupancy unavoidable, and who then
manufactures the arrearage through his own statutory
violations, does not come to this Court with clean hands.
\subsection{\textbf{Ninth Affirmative Defense ---
		Waiver (ORS~\S90.412)}}
To the extent Plaintiff accepted HAP or HRSN payments after
the alleged default without providing written notice of
reservation of rights, Plaintiff waived the right to
terminate for that default. ORS~\S90.412(1). Plaintiff's
own ledger confirms payment receipts throughout the period
in question. \textit{See} Exhibit~B.
% ============================================================
\section{\textbf{DEMAND FOR TRIAL IN THE ALTERNATIVE}}
% ============================================================
Without waiving, and expressly preserving, all
jurisdictional defenses set forth herein, Defendant demands
trial on her counterclaims in the event this Court denies
the Motion to Dismiss or otherwise retains jurisdiction
over any portion of this action. This demand is made solely
to preserve Defendant's right to adjudication of her
compulsory counterclaims under ORS~\S90.370(1)(a) and
ORCP 22~A, and shall not be construed as a general
appearance, a waiver of any jurisdictional objection, or a
concession that this Court has subject matter jurisdiction
over Plaintiff's complaint. ORCP 21~G(1), G(4).
% ============================================================
\section{\textbf{COUNTERCLAIMS}}
% ============================================================
\indent \textbf{Jurisdiction and Mandatory Adjudication.}
ORS~\S90.370(1)(a) provides that in an action for
possession based upon nonpayment of rent, the tenant
\textit{may} counterclaim for any amount recoverable under
the rental agreement or ORS Chapter~90, provided the
landlord had or reasonably should have had knowledge of the
facts constituting the counterclaim prior to filing.
ORS~\S90.370(1)(a) expressly confers jurisdiction on this
Court to hear and determine Defendant's counterclaims in
this FED proceeding, notwithstanding the limited
jurisdiction of the court. Upon the tenant's assertion of
a counterclaim, ORS~\S90.370(1)(b) requires the court to
\textit{determine the amount due to each party} and enter
judgment accordingly. This determination is mandatory, not
discretionary. ORCP 22~A; ORS~\S105.130(4). Each
counterclaim below arises from the same tenancy and is
grounded in Plaintiff's own ledger, text messages, and
written communications, each of which was in Plaintiff's
possession before this action was filed.
\textit{See} Exhibits A--N. Plaintiff had actual and
constructive knowledge of every fact asserted herein.\\
\indent \textbf{Correct Party.} For purposes of these
counterclaims, Plaintiff'' refers to \textbf{JWP WONG
	REAL ESTATE}, the landlord of record and real party in
interest. To the extent the Court finds Jim Wong
individually to be the proper plaintiff, these
counterclaims are asserted against him in the alternative.
Defendant requests that the Court enter judgment on these
counterclaims against the correct party --- JWP WONG REAL
ESTATE --- and grant Defendant equitable, injunctive, and
legal relief accordingly. The real-party defect cannot be
invoked by Plaintiff to defeat compulsory counterclaims.
ORCP 22~A.\\
\indent \textbf{Statute of Limitations.} Claims arising
under ORS Chapter~90 or the rental agreement are subject
to the one-year limitations period of ORS~\S12.125. All
Chapter~90 claims within one year of this filing are timely
on their face. Claims beyond one year are timely under the
continuing tort doctrine, each violation having accrued
anew with each occurrence. The vehicle tow and personal
property conversion claims are governed by ORS~\S12.080
(six-year limitations period for conversion of personal
property), expiring March~6, 2030. All claims are timely
under at least one applicable limitations period.\\
\indent \textbf{Preservation in Event of Leave to Amend.}
In the event the Court finds any counterclaim
insufficiently pleaded, Defendant respectfully requests
leave to amend pursuant to ORCP 23~A to cure any
deficiency. Leave to amend should be freely granted when
justice requires. ORCP 23~A.
\CenteredUnderlined{COUNTERCLAIM DAMAGES SUMMARY}
\begin{longtable}{p{3.3in} p{1.6in} p{.7in}}
	\toprule
	\textbf{Claim} & \textbf{Statutory Basis}
	& \textbf{Amount} \\
	\midrule
	\endhead
	\textbf{Habitability rent abatement.}
	11 documented defects; all noticed in writing;
	Plaintiff failed to remedy after notice; 10\%
	$\times$ \$2,535/mo $\times$ 32~months
	(conservative; court may find higher percentage
	warranted given 28-month drain fly infestation
	and live electrical hazard).
	& ORS \S\S90.320, 90.360(2), 90.370
	& \$8,112.00 \\[12pt]
	\textbf{Renting dwelling known to be unlawful to
		occupy.} Plaintiff knew or reasonably should have
	known the unit was unlawful to occupy: water
	intrusion into live electrical outlet (21 months
	unresolved); 28-month drain fly infestation from
	plumbing defect; standing water and mold from
	leaking master toilet. Remedy is greater of two
	months' periodic rent (\$5,070.00) or twice actual
	habitability damages; floor stated; actual damages
	to be proven at trial.
	& ORS \S90.380
	& \$5,070.00+ \\[12pt]
	\textbf{Unlawful HOA entries into exclusive-use
		patio, backyard, and garage; breach of quiet
		enjoyment.} One month's periodic rent per
	occurrence under ORS~\S90.322(8); 8
	ledger-documented occurrences (Jul.~2024--Apr.~2025);
	Plaintiff confirmed monthly compliance officer
	visits (Exhibit~C, Jul.~30,~2025); pre-ledger
	entry Sep.~25,~2023 establishes additional
	occurrence; each entry independently breaches
	ORS~\S90.360(1); \$2,535 $\times$ 8 $=$
	\$20,280.00 minimum.
	& ORS \S\S90.322(8), 90.360(1), 90.325
	& \$20,280.00 \\[12pt]
	\textbf{Unlawful HOA fee pass-throughs.}
	ORS~\S90.302(1) prohibits the charge;
	ORS~\S90.302(3) provides the remedy: two months'
	periodic rent (\$5,070.00) plus actual damages
	(\$2,220.00 per Exhibit~B) per claim; the demand
	itself is the violation; any lease clause
	authorizing pass-throughs is void under
	ORS~\S90.135.
	& ORS \S\S90.302(1), (3); 90.135
	& \$7,290.00 \\[12pt]
	\textbf{Vehicle tow and conversion --- Dodge
		Journey.} \$14,039.88 outstanding payoff balance
	reported to Defendant's credit (Exhibit~P); towed
	March~6, 2024; no parking pass issued for 8+
	months despite written requests; no valid tow
	authorization (ORS~\S98.854(2)); no pre-tow
	warning; Plaintiff abdicated duty in writing:
	there is nothing I could do''
	(Exhibit~C, Mar.~17,~2024); 6-year SOL under
	ORS~\S12.080 (expires Mar.~6,~2030).
	& ORS \S\S90.325, 90.365; ORS \S12.080
	& \$14,039.88 \\[12pt]
	\textbf{Personal property lost in towed vehicle.}
	iPad, prescription eyeglasses, prescription
	medications, legal and financial paperwork,
	clothing, child's car seat, jewelry, work
	equipment; 6-year SOL under ORS~\S12.080
	(expires Mar.~6,~2030); to be proven at trial.
	& ORS \S\S90.325, 90.365; ORS \S12.080
	& \$3,000.00 \\[12pt]
	\textbf{Retaliation --- 4 independent bases.}
	(1)~Selective HCV disclosure concealing 11 HQS
	violations while forwarding HOA noncompliance
	notice caused by Plaintiff's own failure to
	install window screens;
	(2)~Selective grill enforcement contrary to
	express lease term while similarly situated
	residents are not comparably enforced against
	(Exhibit~N);
	(3)~Bad-faith lease coercion Jan.~25,~2025
	(Exhibit~C);
	(4)~Maintenance of void prior FED judgment
	trapping Defendant in substandard unit;
	statutory remedy up to 3~months' periodic rent
	(\$7,605.00) plus actual damages; to be proven
	at trial.
	& ORS \S90.385(2)
	& TBD \\[12pt]
	\textbf{Emotional distress and consequential
		damages.} Loss of vehicle containing child's car
	seat and prescription medications; credit
	destruction (\$14,039.88 reported to credit,
	Exhibit~P); loss of transportation; impaired
	employment mobility; inability to coordinate
	rental assistance; 32-month habitability pattern;
	HCV voucher threatened by bad-faith disclosure;
	to be proven at trial.
	& ORS \S\S90.365, 90.385; ORS \S90.125
	& TBD \\[12pt]
	\midrule
	\textbf{Minimum Quantified Total
		(before TBD claims)}
	& & \textbf{\$57,791.88} \\[6pt]
	\textit{Plus pre-ledger entry Sep.~25,~2023
		(+\$2,535.00) raises floor to \$60,326.88;
		ORS~\S90.380 actual damages may further increase
		quantified total.}
	& & \textit{(\$60,326.88+)} \\[6pt]
	Plaintiff's corrected rent balance (HOA fines
	and late fees removed; HRSN correctly applied;
	\$0 tenant portion Nov.~2025--Jan.~2026 credited)
	& ORS \S\S90.394, 90.220(9)
	& \$1,345.00 \\[6pt]
	\midrule
	\textbf{Net amount owed to Defendant
		(minimum quantified, before TBD claims)}
	& & \textbf{(\$56,446.88+)} \\
	\bottomrule
\end{longtable}
\indent Defendant's minimum quantified counterclaim damages
of \$57,791.88 exceed the corrected rent balance, resulting
in a net judgment in Defendant's favor of not less than
\$56,446.88. Defendant has satisfied her duty to mitigate
through documented habitability complaints, repeated
attempts to secure alternative housing (frustrated by
Plaintiff's unlawful prior FED judgment), and active
coordination of rental assistance payments.\\
\indent \textbf{A Note on the Magnitude of These
	Counterclaims.} The amounts reflected herein are not
a litigation strategy---they are the product of 32 months
of documented, unresolved habitability failures, repeated
statutory violations, and a manufactured rent arrearage
built from charges the Legislature has expressly prohibited,
including [specify: e.g., unlawful late fees, prohibited
utility charges]. Each dollar figure derives from one of
two sources: either Plaintiff's own ledger, or the remedy
the Oregon Legislature has assigned by statute to the
specific violation Defendant endured. Where the Legislature
has determined that an unlawful entry into a tenant's
exclusive-use space costs one month's periodic rent per
occurrence, or that an unlawful fee demand costs two months'
periodic rent, those are not Defendant's numbers---they are
the Legislature's. Defendant has lived in this home
throughout. These counterclaims are a record of that time.
\indent \textbf{Preservation in Event of Jurisdictional
	Dismissal.} In the event the Court grants Defendant's
Motion to Dismiss, Defendant requests that the Court
either: (a)~retain jurisdiction over these counterclaims
as independent claims pursuant to ORS~\S90.370(1)(a) and
ORCP 22~B, severing them from Plaintiff's dismissed
complaint; or (b)~expressly state in its order of
dismissal that Defendant's counterclaims are dismissed
without prejudice and preserved for assertion in any
subsequent proceeding, including any refiled FED action or
independent civil action, so that no preclusion argument
may be raised against them. Defendant does not abandon
these claims by seeking dismissal of Plaintiff's complaint.
% ============================================================
\section{\textbf{REQUEST FOR RELIEF}}
% ============================================================
WHEREFORE, Defendant Ashle' Penn respectfully requests that
this Court:
\begin{enumerate}[label=\arabic*.,leftmargin=.5in,nosep]
	\item In the alternative, and without waiving any
	jurisdictional defense, grant Defendant a trial on her
	counterclaims pursuant to ORS~\S90.370(1)(a),
	ORCP 22~A, and ORS~\S105.130(3), solely to preserve
	her right to adjudication of those claims in the event
	this Court retains jurisdiction;
	\item Dismiss Plaintiff's Complaint \textbf{with
		prejudice} under ORS~\S90.394(3) and
	\textit{Hickey v.\ Scott}, 370 Or 97 (2022), the
	notice having overstated the amount owed by
	\$4,573.00;
	\item In the alternative, dismiss under
	ORCP 21~A(1)(g) for failure to state a claim on the
	grounds that: (a)~Jim Wong individually is not the
	real party in interest and cannot state a claim for
	possession; and (b)~the void notice is the
	jurisdictional predicate for the entire action;
	\item In the alternative, dismiss under
	ORS~\S90.395 for failure to cooperate with rental
	assistance programs;
	\item In the alternative, dismiss without prejudice
	under ORCP 26~A for failure to prosecute in the name
	of the real party in interest, with leave to refile
	in the name of JWP WONG REAL ESTATE upon submission
	of proof of agency authority;
	\item In the alternative, dismiss for defective
	service under ORCP~7 and ORS~\S105.135;
	\item In the alternative, dismiss or enter default
	due to Plaintiff's non-appearance at the March~2,
	2026 hearing as required by ORS~\S105.137(3);
	\item In the alternative, vacate the June~5, 2026
	trial date and resolve the pending Motion to Dismiss
	before any trial date is set, pursuant to
	ORCP 21~G(4);
	\item Deny Plaintiff's request for possession;
	\item Enter judgment in favor of Defendant on her
	counterclaims against \textbf{JWP WONG REAL ESTATE}
	(or Jim Wong in the alternative) for equitable,
	injunctive, and legal relief, in an amount not less
	than \$57,791.88, resulting in a net judgment in
	Defendant's favor of not less than \$56,446.88, with
	TBD claims reserved for trial;
	\item Grant injunctive relief prohibiting Plaintiff,
	JWP WONG REAL ESTATE, AMS, and the HOA from further
	unauthorized entry into Defendant's exclusive-use
	patio, backyard, and garage without the 24-hour
	notice required by ORS~\S90.322(1)(f);
	\item Grant injunctive relief prohibiting Plaintiff
	from enforcing any HOA demand that Defendant remove
	her outdoor grill, in contravention of the express
	terms of the rental agreement and Defendant's quiet
	enjoyment rights under ORS~\S90.360(1);
	\item Stay any writ of possession pending resolution
	of Defendant's Motion to Dismiss and, if denied,
	pending any appeal, to prevent irreparable harm to
	Defendant and her minor child;
	\item Issue written findings of fact and conclusions
	of law on each jurisdictional and merits issue raised
	herein, including: (a)~whether the notice accurately
	states the amount of rent owed; (b)~whether Jim Wong
	individually is the real party in interest;
	(c)~whether service was legally sufficient;
	(d)~whether Plaintiff appeared at the March~2, 2026
	hearing within the meaning of ORS~\S105.137(3);
	(e)~the amount due to each party on Defendant's
	counterclaims under ORS~\S90.370(1)(b); and
	(f)~the disposition of each counterclaim, for
	purposes of appellate preservation;
	\item Expressly preserve Defendant's counterclaims
	without prejudice in the event of jurisdictional
	dismissal, pursuant to ORCP 22~B and
	ORS~\S90.370(1)(a);
	\item Grant leave to amend any counterclaim found
	insufficiently pleaded, pursuant to ORCP 23~A;
	\item Award Defendant attorney fees and costs under
	ORS~\S90.255; and
	\item Grant such other equitable, injunctive, and
	legal relief as the Court deems just.
\end{enumerate}
\indent I declare under penalty of perjury under the laws
of the State of Oregon that the foregoing is true and
correct to the best of my knowledge and belief.
\vspace{1em}
\noindent Respectfully submitted,\\[4pt]
Date: \underline{March 15, 2026}\\[6pt]
\noindent/s/\includegraphics[scale=0.23]{signature2.png}\\[4pt]
\noindent\rule{3in}{0.5pt}\\
Ashle' Penn, Defendant, \textit{In Propria Persona}\\
17548 NW Springville Rd \#9, Portland, OR 97229\\[4pt]
\textit{Defendant appears pro se and specially,
	without waiving any jurisdictional defenses.}
% ============================================================
% EXHIBIT LIST
% ============================================================
\newpage
\fancyfoot[L]{\fontsize{8}{10}\selectfont
	\textbf{\textsc{EXHIBIT LIST}} --- Pg.~\thepage}
\begin{center}
	{\normalsize\textbf{EXHIBIT LIST}}\\[4pt]
	{\normalsize\textit{Jim Wong v.\ Ashle' Penn},
		Case No.\ 26LT03815}\\[4pt]
	{\normalsize Washington County Circuit Court}
\end{center}
\vspace{0.5em}
\begin{longtable}{p{0.9in} p{4.7in}}
	\toprule
	\textbf{Exhibit} & \textbf{Description} \\
	\midrule
	\endhead
	Exhibit A & Defendant's Counterclaim Damages Summary
	(Section~VIII of this Amended Answer). \\[6pt]
	Exhibit B & Plaintiff's Rent Ledger, 17548 NW
	Springville Rd Unit~9, Tenant Ashle' Penn,
	July~2024--January~2026, transmitted by email from
	jimwong@me.com on January~21, 2026 at 11:56~p.m.;
	confirming HOA fines of \$2,220.00, late fees of
	\$1,045.00, and HRSN credit of \$4,494.00 applied
	to November~2025 only; basis for Plaintiff's claimed
	balance of \$5,918.00. \\[6pt]
	Exhibit C & Text Message and Email Record, Jim Wong
	and Ashle' Penn, July~2023--March~2026, including:
	habitability complaints and responses; HOA
	noncompliance notices forwarded by Plaintiff;
	Plaintiff's July~30, 2025 admission of monthly
	compliance officer visits; Defendant's January~13,
	2025 explanation of window tape as self-help remedy
	for missing screens; Plaintiff's March~11--18, 2024
	tow communications and admissions including there
	is nothing I could do''; failure to provide parking
	pass; failure to provide replacement mail key;
	January~25, 2025 lease coercion; Plaintiff's
	January~21, 2026 email admitting HOA fine suspension
	was discretionary; Plaintiff's March~2, 2026 text
	message at 10:19~a.m.\ confirming ex parte contact
	with the court regarding Defendant's pending motion
	before Defendant was contacted. \\[6pt]
	Exhibit D & Lease Agreement, 17548 NW Springville
	Rd Unit~9, JWP WONG REAL ESTATE and Ashle' Penn,
	including provision expressly permitting outdoor
	grill. \\[6pt]
	Exhibit E & CC\&Rs for Courtyards at Springville,
	all phases: exclusive-use limited common element
	designations (patios, backyards, garages);
	Sections~5.1 and~6.1. \\[6pt]
	Exhibit F & Bylaws, Courtyards at Springville:
	Sections~5.3(a), 5.6, 5.7 confirming HOA
	assessments and enforcement remedies run against
	unit owners only. \\[6pt]
	Exhibit G & WCHA HAP Payment History, Ashle' Penn,
	17548 NW Springville Rd Unit~9, confirming payments
	to JWP WONG REAL ESTATE; \$0.00 tenant portion for
	November~2025--January~2026 following income
	recertification. \\[6pt]
	Exhibit H & Income Recertification Signature
	Document, Ashle' Penn, confirming \$0.00 tenant
	portion effective November~2025. \\[6pt]
	Exhibit I & HOA Compliance Letter, Case No.\
	2024-CSVC-00058, August~13, 2024. \\[6pt]
	Exhibit J & HOA Compliance Letter, Case No.\
	2024-CSVC-00058, September~17, 2024. \\[6pt]
	Exhibit K & HOA Compliance Letter, Case No.\
	2024-CSVC-00059, September~24, 2024. \\[6pt]
	Exhibit L & HOA Compliance Letter, Case No.\
	2024-CSVC-00059, January~13, 2025. \\[6pt]
	Exhibit M & Text message screenshot, March~2, 2026
	at 10:19~a.m., Jim Wong to Ashle' Penn, confirming
	Plaintiff's ex parte contact with the court
	regarding Defendant's pending motion before
	Defendant was contacted. \\[6pt]
	Exhibit N & Photograph of parking area at 17548 NW
	Springville Rd showing absence of fire lane markings
	at location of tow; and photograph showing outdoor
	grills maintained by other residents without
	comparable enforcement action. \\[6pt]
	Exhibit O & Jim Wong email(s), March~17--18, 2024,
	admitting: no pre-tow warning was given; Defendant's
	photographs showed vehicle was not in the fire lane;
	and there is nothing I could do'' --- establishing
	Plaintiff's abdication of his non-delegable duty
	under ORS~\S90.325. \\[6pt]
	Exhibit P & Defendant's credit report and/or bank
	statement reflecting \$14,039.88 outstanding payoff
	balance on Dodge Journey reported to credit
	following tow. \\
	\bottomrule
\end{longtable}
% ============================================================
% EXHIBIT A — COUNTERCLAIM DAMAGES SUMMARY (cover sheet)
% ============================================================
\newpage
\fancyfoot[L]{\fontsize{8}{10}\selectfont
	\textbf{\textsc{DEFENDANT'S EXHIBIT A}} ---
	Pg.~\thepage}
\begin{flushleft}
	\textbf{\large Defendant's Exhibit A}
\end{flushleft}
\begin{center}
	\textbf{COUNTERCLAIM DAMAGES SUMMARY}\\[4pt]
	\textit{Jim Wong v.\ Ashle' Penn},
	Case No.\ 26LT03815\\[4pt]
	Washington County Circuit Court
\end{center}
\vspace{0.5em}
\noindent \textit{See} Section~VIII of Defendant's Amended
Answer, incorporated herein by reference. The counterclaim
damages summary set forth in Section~VIII constitutes
Defendant's Exhibit~A.
% ============================================================
% EXHIBIT B — LEDGER
% ============================================================
\newpage
\fancyfoot[L]{\fontsize{8}{10}\selectfont
	\textbf{\textsc{DEFENDANT'S EXHIBIT B}} ---
	Pg.~\thepage}
\begin{flushleft}
	\textbf{\large Defendant's Exhibit B}
\end{flushleft}
\begin{center}
	\textbf{PLAINTIFF'S RENT LEDGER --- UNIT 9}\\[4pt]
	\textit{Jim Wong v.\ Ashle' Penn},
	Case No.\ 26LT03815\\[4pt]
	Transmitted January~21, 2026 at 11:56~p.m.\
	by jimwong@me.com
\end{center}
\vspace{0.5em}
\includepdf[pages=-,pagecommand={}]{EXHIBIT.Ashle F9 Ledger - Dec 2026.pdf}
% ============================================================
% EXHIBIT C — TEXT / EMAIL RECORD
% ============================================================
\newpage
\fancyfoot[L]{\fontsize{8}{10}\selectfont
	\textbf{\textsc{DEFENDANT'S EXHIBIT C}} ---
	Pg.~\thepage}
\begin{flushleft}
	\textbf{\large Defendant's Exhibit C}
\end{flushleft}
\begin{center}
	\textbf{TEXT MESSAGE AND EMAIL RECORD ---
		SELECTED RELEVANT COMMUNICATIONS}\\[4pt]
	\textit{Jim Wong v.\ Ashle' Penn},
	Case No.\ 26LT03815
\end{center}
\vspace{0.5em}
\noindent \textbf{Index of Communications:}
\begin{longtable}{p{1.2in} p{4.4in}}
	\toprule
	\textbf{Date} & \textbf{Subject} \\
	\midrule
	\endhead
	Jul.~2023 & Defendant's vehicle registration and
	parking pass request \\[3pt]
	Nov.~16, 2023 & Drain fly infestation --- written
	notice to Plaintiff \\[3pt]
	Jun.~26, 2024 & Leaking kitchen sink; water
	intrusion into live electrical outlet ---
	written notice \\[3pt]
	Mar.~11--18, 2024 & Tow communications: parking
	pass failure; shifting justifications; Plaintiff's
	admission there is nothing I could do'' \\[3pt]
	Jan.~13, 2025 & Window tape as self-help remedy
	for missing screens \\[3pt]
	Jan.~25, 2025 & Bad-faith lease coercion \\[3pt]
	Jul.~30, 2025 & Plaintiff's admission: They will
	send compliance officer to check on whether they
	have been addressed and accumulate the fines each
	month'' \\[3pt]
	Nov.~17, 2025 & Leaking master toilet; standing
	water; mold --- written notice \\[3pt]
	Nov.~2025--Jan.~2026 & Non-functioning stove
	burner; inadequate hot water; power outages;
	non-functioning outlets; mail key requests \\[3pt]
	Jan.~21, 2026 & Plaintiff's email transmitting
	ledger; admission HOA fine suspension was
	discretionary: the HOA kindly agreed to
	temporarily suspend'' \\[3pt]
	Mar.~2, 2026 (10:19~a.m.) & Plaintiff's text
	confirming ex parte court contact: I spoke with
	the court regarding your motion to file
	documents'' \\
	\bottomrule
\end{longtable}
\vspace{0.5em}
\includegraphics[width=\textwidth,keepaspectratio]{sent.png}
% ============================================================
% EXHIBIT D — LEASE
% ============================================================
\newpage
\fancyfoot[L]{\fontsize{8}{10}\selectfont
	\textbf{\textsc{DEFENDANT'S EXHIBIT D}} ---
	Pg.~\thepage}
\begin{flushleft}
	\textbf{\large Defendant's Exhibit D}
\end{flushleft}
\begin{center}
	\textbf{LEASE AGREEMENT ---
		17548 NW SPRINGVILLE RD UNIT 9}\\[4pt]
	\textit{Jim Wong v.\ Ashle' Penn},
	Case No.\ 26LT03815
\end{center}
\vspace{0.5em}
\includepdf[pages=-,pagecommand={}]{17548 LEASE.EXHIBIT - Copy.pdf}
% ============================================================
% EXHIBIT E — CC&Rs
% ============================================================
\newpage
\fancyfoot[L]{\fontsize{8}{10}\selectfont
	\textbf{\textsc{DEFENDANT'S EXHIBIT E}} ---
	Pg.~\thepage}
\begin{flushleft}
	\textbf{\large Defendant's Exhibit E}
\end{flushleft}
\begin{center}
	\textbf{CC\&Rs --- COURTYARDS AT SPRINGVILLE,
		ALL PHASES}\\[4pt]
	\textit{Jim Wong v.\ Ashle' Penn},
	Case No.\ 26LT03815
\end{center}
\vspace{0.5em}
\includepdf[pages=-,pagecommand={}]{EXHIBIT. CCRS All Phases.pdf}
% ============================================================
% EXHIBIT F — BYLAWS
% ============================================================
\newpage
\fancyfoot[L]{\fontsize{8}{10}\selectfont
	\textbf{\textsc{DEFENDANT'S EXHIBIT F}} ---
	Pg.~\thepage}
\begin{flushleft}
	\textbf{\large Defendant's Exhibit F}
\end{flushleft}
\begin{center}
	\textbf{BYLAWS --- COURTYARDS AT SPRINGVILLE}\\[4pt]
	\textit{Jim Wong v.\ Ashle' Penn},
	Case No.\ 26LT03815
\end{center}
\vspace{0.5em}
\includepdf[pages=-,pagecommand={}]{EXHIBIT. Bylaws (1) - Copy.pdf}
% ============================================================
% EXHIBIT G — HAP PAYMENT HISTORY
% ============================================================
\newpage
\fancyfoot[L]{\fontsize{8}{10}\selectfont
	\textbf{\textsc{DEFENDANT'S EXHIBIT G}} ---
	Pg.~\thepage}
\begin{flushleft}
	\textbf{\large Defendant's Exhibit G}
\end{flushleft}
\begin{center}
	\textbf{WCHA HAP PAYMENT HISTORY}\\[4pt]
	\textit{Jim Wong v.\ Ashle' Penn},
	Case No.\ 26LT03815
\end{center}
\vspace{0.5em}
\includepdf[pages=-,pagecommand={}]{EXHIBIT. TenantHAP_2-28-2026_4-42.pdf}
% ============================================================
% EXHIBIT H — RECERTIFICATION
% ============================================================
\newpage
\fancyfoot[L]{\fontsize{8}{10}\selectfont
	\textbf{\textsc{DEFENDANT'S EXHIBIT H}} ---
	Pg.~\thepage}
\begin{flushleft}
	\textbf{\large Defendant's Exhibit H}
\end{flushleft}
\begin{center}
	\textbf{INCOME RECERTIFICATION SIGNATURE
		DOCUMENT}\\[4pt]
	\textit{Jim Wong v.\ Ashle' Penn},
	Case No.\ 26LT03815
\end{center}
\vspace{0.5em}
\includepdf[pages=-,pagecommand={}]{EXHIBIT. Recert__Ashle'_Penn_44181362_Member_2309917 Signature Document.pdf}
% ============================================================
% EXHIBIT I — COMPLIANCE LETTER 2024-CSVC-00058, 8.13.24
% ============================================================
\newpage
\fancyfoot[L]{\fontsize{8}{10}\selectfont
	\textbf{\textsc{DEFENDANT'S EXHIBIT I}} ---
	Pg.~\thepage}
\begin{flushleft}
	\textbf{\large Defendant's Exhibit I}
\end{flushleft}
\begin{center}
	\textbf{HOA COMPLIANCE LETTER ---
		CASE NO.\ 2024-CSVC-00058}\\[2pt]
	\textbf{August~13, 2024}\\[4pt]
	\textit{Jim Wong v.\ Ashle' Penn},
	Case No.\ 26LT03815
\end{center}
\vspace{0.5em}
\includepdf[pages=-,pagecommand={}]{EXHIBIT.F9 Compliance Letter 8.13.24.pdf}
% ============================================================
% EXHIBIT J — COMPLIANCE LETTER 2024-CSVC-00058, 9.17.24
% ============================================================
\newpage
\fancyfoot[L]{\fontsize{8}{10}\selectfont
	\textbf{\textsc{DEFENDANT'S EXHIBIT J}} ---
	Pg.~\thepage}
\begin{flushleft}
	\textbf{\large Defendant's Exhibit J}
\end{flushleft}
\begin{center}
	\textbf{HOA COMPLIANCE LETTER ---
		CASE NO.\ 2024-CSVC-00058}\\[2pt]
	\textbf{September~17, 2024}\\[4pt]
	\textit{Jim Wong v.\ Ashle' Penn},
	Case No.\ 26LT03815
\end{center}
\vspace{0.5em}
\includepdf[pages=-,pagecommand={}]{EXHIBIT_ComplianceLetter_00058_9-17-24.pdf}
% ============================================================
% EXHIBIT K — COMPLIANCE LETTER 2024-CSVC-00059, 9.24.24
% ============================================================
\newpage
\fancyfoot[L]{\fontsize{8}{10}\selectfont
	\textbf{\textsc{DEFENDANT'S EXHIBIT K}} ---
	Pg.~\thepage}
\begin{flushleft}
	\textbf{\large Defendant's Exhibit K}
\end{flushleft}
\begin{center}
	\textbf{HOA COMPLIANCE LETTER ---
		CASE NO.\ 2024-CSVC-00059}\\[2pt]
	\textbf{September~24, 2024}\\[4pt]
	\textit{Jim Wong v.\ Ashle' Penn},
	Case No.\ 26LT03815
\end{center}
\vspace{0.5em}
\includepdf[pages=-,pagecommand={}]{EXHIBIT.Compliance Letter 9.24.24.pdf}
% ============================================================
% EXHIBIT L — COMPLIANCE LETTER 2024-CSVC-00059, 1.13.25
% ============================================================
\newpage
\fancyfoot[L]{\fontsize{8}{10}\selectfont
	\textbf{\textsc{DEFENDANT'S EXHIBIT L}} ---
	Pg.~\thepage}
\begin{flushleft}
	\textbf{\large Defendant's Exhibit L}
\end{flushleft}
\begin{center}
	\textbf{HOA COMPLIANCE LETTER ---
		CASE NO.\ 2024-CSVC-00059}\\[2pt]
	\textbf{January~13, 2025}\\[4pt]
	\textit{Jim Wong v.\ Ashle' Penn},
	Case No.\ 26LT03815
\end{center}
\vspace{0.5em}
\includepdf[pages=-,pagecommand={}]{EXHIBIT_ComplianceLetter_00059_1-13-25.pdf}
% ============================================================
% EXHIBIT M — EX PARTE TEXT
% ============================================================
\newpage
\fancyfoot[L]{\fontsize{8}{10}\selectfont
	\textbf{\textsc{DEFENDANT'S EXHIBIT M}} ---
	Pg.~\thepage}
\begin{flushleft}
	\textbf{\large Defendant's Exhibit M}
\end{flushleft}
\begin{center}
	\textbf{TEXT MESSAGE --- PLAINTIFF'S EX PARTE
		COURT CONTACT}\\[2pt]
	\textbf{March~2, 2026 at 10:19~a.m.}\\[4pt]
	\textit{Jim Wong v.\ Ashle' Penn},
	Case No.\ 26LT03815
\end{center}
\vspace{0.5em}
\includegraphics[width=\textwidth,keepaspectratio]{sent.png}
% ============================================================
% EXHIBIT N — PARKING / GRILL PHOTOGRAPHS
% ============================================================
\newpage
\fancyfoot[L]{\fontsize{8}{10}\selectfont
	\textbf{\textsc{DEFENDANT'S EXHIBIT N}} ---
	Pg.~\thepage}
\begin{flushleft}
	\textbf{\large Defendant's Exhibit N}
\end{flushleft}
\begin{center}
	\textbf{PHOTOGRAPHS --- PARKING AREA AND
		COMPARATIVE GRILL ENFORCEMENT}\\[4pt]
	\textit{Jim Wong v.\ Ashle' Penn},
	Case No.\ 26LT03815
\end{center}
\vspace{0.5em}
\noindent \textbf{Figure 1:} Parking area at 17548 NW
Springville Rd showing absence of fire lane markings
at location of tow, March~6, 2024.
\vspace{0.5em}
\includegraphics[width=\textwidth,keepaspectratio]{EXHIBIT. redzone1.jpg}
% ============================================================
% EXHIBIT O — ABDICATION EMAIL
% ============================================================
\newpage
\fancyfoot[L]{\fontsize{8}{10}\selectfont
	\textbf{\textsc{DEFENDANT'S EXHIBIT O}} ---
	Pg.~\thepage}
\begin{flushleft}
	\textbf{\large Defendant's Exhibit O}
\end{flushleft}
\begin{center}
	\textbf{JIM WONG EMAIL --- ABDICATION OF
		LANDLORD DUTY}\\[2pt]
	\textbf{March~17--18, 2024}\\[4pt]
	\textit{Jim Wong v.\ Ashle' Penn},
	Case No.\ 26LT03815
\end{center}
\vspace{0.5em}
\noindent This exhibit contains Plaintiff's written
communications of March~17--18, 2024, in which Plaintiff
acknowledged that: (1)~a warning was required before
towing; (2)~Defendant's photographs showed her vehicle
was not in the fire lane; and (3)~there is nothing I
could do'' --- establishing Plaintiff's abdication of his
non-delegable duty under ORS~\S90.325.
\vspace{0.5em}
\includepdf[pages=-,pagecommand={}]{EXHIBIT.font of car in front of garage.jpg}
% ============================================================
% EXHIBIT P — CREDIT / VEHICLE PAYOFF
% ============================================================
\newpage
\fancyfoot[L]{\fontsize{8}{10}\selectfont
	\textbf{\textsc{DEFENDANT'S EXHIBIT P}} ---
	Pg.~\thepage}
\begin{flushleft}
	\textbf{\large Defendant's Exhibit P}
\end{flushleft}
\begin{center}
	\textbf{CREDIT REPORT / VEHICLE PAYOFF BALANCE}\\[4pt]
	\textit{Jim Wong v.\ Ashle' Penn},
	Case No.\ 26LT03815
\end{center}
\vspace{0.5em}
\noindent This exhibit contains Defendant's credit report
and/or bank statement reflecting the \$14,039.88
outstanding payoff balance on the Dodge Journey reported
to Defendant's credit following the March~6, 2024 tow.
\vspace{1em}
\noindent \textit{[Attach credit report / bank statement
	here before filing.]}
% ============================================================
% DOCUMENT 2 OF 4 — NOTICE OF MOTION
% ============================================================
\newpage
\fancyfoot[L]{\fontsize{8}{10}\selectfont
	\textbf{\textsc{NOTICE OF MOTION}} ---
	Pg.~\thepage}
\begin{center}
	{\normalsize\textbf{IN THE CIRCUIT COURT OF THE
			STATE OF OREGON}}\\[2pt]
	{\normalsize\textbf{FOR THE COUNTY OF
			WASHINGTON}}\\[2pt]
\end{center}
\CaptionBlock{\textbf{26LT03815}}{%
	\begin{center}
		\fontsize{9}{9}\selectfont
		\textbf{NOTICE OF DEFENDANT'S RENEWED\\
			MOTION TO DISMISS FOR LACK OF\\
			SUBJECT MATTER JURISDICTION\\
			AND LACK OF STANDING}%
	\end{center}
}{%
	JIM WONG,\\[4pt]
	\hspace{0.3in}Plaintiff,\\[8pt]
	vs.\\[8pt]
	ASHLE' PENN,\\[8pt]
	\hspace{0.3in}Defendant,
	\textit{In Propria Persona}.%
}
\vspace{1em}
PLEASE TAKE NOTICE that Defendant Ashle' Penn, appearing
specially and without waiving any jurisdictional defense,
hereby renews her Motion to Dismiss for Lack of Subject
Matter Jurisdiction and Lack of Standing, as set forth in
her Amended Answer filed herewith.
\indent This motion is based upon the following grounds,
each of which is independently sufficient to require
dismissal:
\begin{enumerate}[label=\arabic*.,leftmargin=.5in,
	itemsep=2pt]
	\item \textbf{Defective Notice of Termination} ---
	Plaintiff's notice overstates the amount of rent owed
	by \$4,573.00 (noticed: \$5,918.00; corrected:
	\$1,345.00), rendering the notice void and depriving
	this Court of subject matter jurisdiction under
	ORS~\S90.394(3) and \textit{Hickey v.\ Scott},
	370 Or 97 (2022);
	\item \textbf{Wrong Plaintiff; Failure to State a
		Claim} --- Jim Wong individually is not the real
	party in interest; the landlord of record is JWP WONG
	REAL ESTATE, LLC; a plaintiff without a cognizable
	legal right to possession cannot state a claim under
	ORCP 21~A(1)(g) and ORS~\S105.130(4)
%% PICKS UP from: ORS~\S105.130(4
\item \textbf{Indispensable Party Not Joined} ---
JWP WONG REAL ESTATE is an indispensable party whose
absence deprives this Court of the ability to grant
complete relief; ORCP 26~A; ORCP 21~A(1)(g);
\item \textbf{Defective Service} --- Service did not
comply with ORCP~7 and ORS~\S105.135, depriving this
Court of personal jurisdiction over Defendant;
ORCP 21~A(1)(e); ORCP 21~G(1);
\item \textbf{Failure to Cooperate with Rental
	Assistance} --- Plaintiff accepted HAP and HRSN
payments and misapplied them to unlawful charges in
violation of ORS~\S90.220(9); ORS~\S90.395
independently requires dismissal;
\item \textbf{Plaintiff's Non-Appearance} ---
Plaintiff failed to appear when this case was called
on March~2, 2026; ORS~\S105.137(3) mandates
dismissal or entry of judgment for Defendant upon
Plaintiff's failure to appear; and
\item \textbf{Trial Date Set Over Pending
	Jurisdictional Motion} --- No trial date may be set
while a jurisdictional motion remains pending and
unresolved; ORCP 21~G(4); ORS~\S105.130(4).
\end{enumerate}
\indent This motion is supported by Defendant's Amended
Answer, Affirmative Defenses, and Counterclaims filed
herewith, including all exhibits thereto, each of which
is incorporated herein by reference.
\indent Defendant respectfully requests that this Court
resolve this motion before any trial date proceeds and
before requiring further substantive filings from
Defendant. ORCP 21~G(4).
\vspace{1em}
\noindent Respectfully submitted,\\[4pt]
Date: \underline{March 15, 2026}\\[6pt]
\noindent/s/\includegraphics[scale=0.23]{signature2.png}\\[4pt]
\noindent\rule{3in}{0.5pt}\\
Ashle' Penn, Defendant, \textit{In Propria Persona}\\
17548 NW Springville Rd \#9, Portland, OR 97229\\[4pt]
\textit{Defendant appears pro se and specially,
without waiving any jurisdictional defenses.}
% ============================================================
% DOCUMENT 3 OF 4 — PROPOSED ORDER
% ============================================================
\newpage
\fancyfoot[L]{\fontsize{8}{10}\selectfont
\textbf{\textsc{PROPOSED ORDER}} ---
Pg.~\thepage}
\begin{center}
{\normalsize\textbf{IN THE CIRCUIT COURT OF THE
		STATE OF OREGON}}\\[2pt]
{\normalsize\textbf{FOR THE COUNTY OF
		WASHINGTON}}\\[2pt]
		\end{center}
		\CaptionBlock{\textbf{26LT03815}}{%
\begin{center}
	\fontsize{9}{9}\selectfont
	\textbf{PROPOSED ORDER ON DEFENDANT'S\\
		RENEWED MOTION TO DISMISS}%
\end{center}
}{%
JIM WONG,\\[4pt]
\hspace{0.3in}Plaintiff,\\[8pt]
vs.\\[8pt]
ASHLE' PENN,\\[8pt]
\hspace{0.3in}Defendant,
\textit{In Propria Persona}.%
}
\vspace{1em}
THIS MATTER having come before the Court on Defendant's
Renewed Motion to Dismiss for Lack of Subject Matter
Jurisdiction and Lack of Standing, and the Court having
considered the motion, the Amended Answer, and all
exhibits thereto, and being duly advised in the premises:
\vspace{0.5em}
\noindent \textbf{IT IS HEREBY ORDERED} as follows:
\vspace{0.5em}
\noindent \textbf{[OPTION A --- GRANT MOTION TO DISMISS]}
\vspace{0.3em}
\noindent $\square$ Defendant's Renewed Motion to Dismiss
is \textbf{GRANTED}.
\vspace{0.3em}
\noindent $\square$ Plaintiff's Complaint is dismissed
\underline{\hspace{1in}} with prejudice
\underline{\hspace{1in}} without prejudice, on the
following ground(s):
\vspace{0.3em}
\begin{enumerate}[label=$\square$,leftmargin=.5in,
itemsep=2pt]
\item Defective notice of termination; notice void
under ORS~\S90.394(3) and \textit{Hickey v.\ Scott},
370 Or 97 (2022); this Court lacks subject matter
jurisdiction.
\item Wrong plaintiff; Jim Wong individually is not
the real party in interest; failure to state a claim
under ORCP 21~A(1)(g).
\item Indispensable party JWP WONG REAL ESTATE not
joined; ORCP 26~A.
\item Defective service; ORCP 21~A(1)(e).
\item Failure to cooperate with rental assistance;
ORS~\S90.395.
\item Plaintiff's non-appearance; ORS~\S105.137(3).
\end{enumerate}
\vspace{0.3em}
\noindent $\square$ The June~5, 2026 trial date is
\textbf{VACATED}.
\vspace{0.3em}
\noindent $\square$ Defendant's counterclaims are
dismissed \textbf{WITHOUT PREJUDICE} and are expressly
preserved for assertion in any subsequent proceeding,
including any refiled FED action or independent civil
action. No preclusion shall arise from this dismissal
as to Defendant's counterclaims. ORCP 22~B;
ORS~\S90.370(1)(a).
\vspace{0.3em}
\noindent $\square$ In the alternative, this Court
retains jurisdiction over Defendant's counterclaims as
independent claims pursuant to ORS~\S90.370(1)(a) and
ORCP 22~B, severed from Plaintiff's dismissed complaint.
\vspace{0.5em}
\noindent \textbf{[OPTION B --- DENY MOTION TO DISMISS]}
\vspace{0.3em}
\noindent $\square$ Defendant's Renewed Motion to Dismiss
is \textbf{DENIED}.
\vspace{0.3em}
\noindent $\square$ This matter is set for trial on
\underline{\hspace{1.5in}}, 2026 at
\underline{\hspace{0.75in}} a.m./p.m., before Judge
\underline{\hspace{1.5in}}.
\vspace{0.3em}
\noindent $\square$ All jurisdictional objections raised
by Defendant are \textbf{PRESERVED} for appellate review.
\vspace{0.3em}
\noindent $\square$ The Court shall determine the amount
due to each party on Defendant's counterclaims at trial
pursuant to ORS~\S90.370(1)(b). This determination is
mandatory.
\vspace{0.5em}
\noindent \textbf{[ALL ORDERS]}
\vspace{0.3em}
\noindent $\square$ Any writ of possession is
\textbf{STAYED} pending resolution of this motion and,
if denied, pending any appeal.
\vspace{0.3em}
\noindent $\square$ The Court issues written findings of
fact and conclusions of law on the following:
\begin{enumerate}[label=(\alph*),leftmargin=.5in,
itemsep=2pt]
\item Whether the notice accurately states the
amount of rent owed;
\item Whether Jim Wong individually is the real
party in interest;
\item Whether service was legally sufficient;
\item Whether Plaintiff appeared at the March~2,
2026 hearing within the meaning of
ORS~\S105.137(3);
\item The amount due to each party on Defendant's
counterclaims under ORS~\S90.370(1)(b); and
\item The disposition of each counterclaim.
\end{enumerate}
\vspace{0.3em}
\noindent $\square$ Defendant is awarded attorney fees
and costs under ORS~\S90.255 in the amount of
\$\underline{\hspace{1in}}.
\vspace{1.5em}
\noindent IT IS SO ORDERED.
\vspace{1.5em}
\noindent DATED this \underline{\hspace{0.5in}} day of
\underline{\hspace{1.2in}}, 2026.
\vspace{2em}
\noindent\rule{3.5in}{0.5pt}\\
Circuit Court Judge\\
Washington County Circuit Court
% ============================================================
% DOCUMENT 4 OF 4 — CERTIFICATE OF SERVICE
% ============================================================
\newpage
\fancyfoot[L]{\fontsize{8}{10}\selectfont
\textbf{\textsc{CERTIFICATE OF SERVICE}} ---
Pg.~\thepage}
\begin{center}
{\normalsize\textbf{IN THE CIRCUIT COURT OF THE
		STATE OF OREGON}}\\[2pt]
{\normalsize\textbf{FOR THE COUNTY OF
		WASHINGTON}}\\[2pt]
		\end{center}
		\CaptionBlock{\textbf{26LT03815}}{%
\begin{center}
	\fontsize{9}{9}\selectfont
	\textbf{CERTIFICATE OF SERVICE}%
\end{center}
}{%
JIM WONG,\\[4pt]
\hspace{0.3in}Plaintiff,\\[8pt]
vs.\\[8pt]
ASHLE' PENN,\\[8pt]
\hspace{0.3in}Defendant,
\textit{In Propria Persona}.%
}
\vspace{1em}
I, Ashle' Penn, hereby certify that on
\underline{\hspace{0.5in}} March 2026, I served a true
and correct copy of the following documents:
\begin{enumerate}[label=\arabic*.,leftmargin=.5in,
itemsep=2pt]
\item Defendant's Amended Answer, Affirmative
Defenses, and Renewed Motion to Dismiss, with
Exhibits A through P;
\item Notice of Defendant's Renewed Motion to
Dismiss; and
\item Proposed Order on Defendant's Renewed Motion
to Dismiss
\end{enumerate}
\noindent upon Plaintiff Jim Wong by the following
method:
\vspace{0.5em}
\noindent $\square$ \textbf{First-class U.S.\ mail,}
postage prepaid, addressed to:
\vspace{0.5em}
\noindent \hspace{0.5in} Jim Wong\\
\hspace{0.5in} \underline{\hspace{3in}}\\
\hspace{0.5in} \underline{\hspace{3in}}
\vspace{0.5em}
\noindent $\square$ \textbf{Electronic service} to:
\underline{jimwong@me.com}
\vspace{0.5em}
\noindent $\square$ \textbf{Personal service} at:
\underline{\hspace{3in}}
\vspace{1em}
\noindent I declare under penalty of perjury under the
laws of the State of Oregon that the foregoing is true
and correct.
\vspace{1.5em}
\noindent Date: \underline{\hspace{1.5in}}\\[6pt]
\noindent/s/\includegraphics[scale=0.23]{signature2.png}\\[4pt]
\noindent\rule{3in}{0.5pt}\\
Ashle' Penn, Defendant, \textit{In Propria Persona}\\
17548 NW Springville Rd \#9, Portland, OR 97229
\end{document}
